\documentclass[11pt,letterpaper]{article}

% Essential Packages
\usepackage[utf8]{inputenc}
\usepackage{amsmath}
\usepackage{amssymb}
\usepackage{amsfonts}
\usepackage{graphicx}
\usepackage{hyperref}
\usepackage{geometry}
\usepackage{cite}
\usepackage{xcolor}

% Margin Configuration
\geometry{margin=1in}

% Hyperlink Configuration
\hypersetup{
    colorlinks=true,
    linkcolor=blue,
    filecolor=magenta,      
    urlcolor=blue,
    citecolor=blue,
    pdftitle={A Mechanical Extension of General Relativity: Redefining the Singularity as a State of Equilibrium},
    pdfauthor={Robin Skavberg}
}

\title{\textbf{A Mechanical Extension of General Relativity: Redefining the Singularity as a State of Equilibrium}}

\author{
    \textbf{Robin Skavberg} \\
    \small Independent Researcher \\
    \small ORCID: \href{https://orcid.org/0009-0003-9126-6196}{0009-0003-9126-6196} \\
    \small Email: \href{mailto:skavberg@gmail.com}{skavberg@gmail.com}
}

\date{\today}

\begin{document}

\maketitle

\begin{abstract}
General Relativity (GR) provides an exceptionally robust description of spacetime geometry; however, its predictive power meets a formal limit at the singularity. This paper proposes a structural extension to the Einstein-Hilbert framework by applying the definition of the gravitational constant as a dynamic ratio $G = c^2 / (8\pi + P_m)$ established in prior work \cite{Skavberg2026}. We propose a redefinition of the term singularity, identifying it not as a mathematical point of infinite density, but as a physical mechanical saturation limit. By introducing this boundary condition, we demonstrate that the manifold reaches a structural limit of compression, resulting in the singularity being in a state of equilibrium. This resolves the Black Hole Information Paradox by ensuring that quantum states are preserved within a saturated volumetric structure rather than destroyed at an infinitesimal point.
\end{abstract}

\section{Introduction: Redefining the Singularity}
The mathematical prediction of a singularity in standard General Relativity suggests a breakdown of causality and predictability. Building upon the mechanical derivation of $G$ established in recent literature \cite{Skavberg2026}, this paper proposes that the singularity maintains a state of equilibrium due to mechanical saturation. 

Under the Stiffness Matching Principle (SMP), the singularity represents the threshold where matter-field pressure ($P_m$) matches the intrinsic stiffness of the vacuum ($8\pi$). At this limit, the interaction between matter and geometry saturates, establishing a saturated equilibrium state. This effectively halts gravitational collapse at a finite density. The effective equilibrium allows for a consistent description of black hole interiors that satisfies the requirements of both General Relativity and quantum unitarity.

\section{From Infinity to Saturation}
To redefine the singularity as a state of equilibrium, we apply a cause-and-effect analysis to the structural limit of the spacetime manifold:

\begin{enumerate}
    \item \textbf{Cause (Increased Field Pressure $P_m$):} Gravitational collapse drives the localized energy density $P_m$ toward extreme limits, increasing the stress on the manifold.
    \item \textbf{Mechanism (Stiffness Matching):} As established in the mechanical route to $G$ \cite{Skavberg2026}, the coupling adjusts as the field pressure approaches the intrinsic vacuum stiffness ($8\pi$).
    \item \textbf{Effect (Saturated Equilibrium):} The singularity is reached not as an infinity, but as a point where the manifold can no longer sustain additional curvature. This mechanical saturation creates a steady-state equilibrium, forming a stable core that halts singular collapse and maintains structural integrity.
\end{enumerate}



\section{Formalizing the Dynamic Coupling}
The SMP serves as the analytical tool for this redefinition, allowing $G$ to respond to the tensor strength of the field to maintain equilibrium:
\begin{equation}
G(P_m) = \frac{c^2}{8\pi + P_m}
\end{equation}
The effective coupling constant $\kappa_{eff}$ thus becomes:
\begin{equation}
\kappa_{eff} = \frac{8\pi G}{c^4} = \frac{8\pi}{c^2(8\pi + P_m)}
\end{equation}
In high-pressure regimes, $\kappa_{eff}$ demonstrates the saturation of the manifold, establishing the equilibrium state that replaces mathematical infinity with physical stability.

\section{Information Preservation in the Saturated Equilibrium}
By redefining the singularity as a state of equilibrium, we resolve the Black Hole Information Paradox:
\begin{itemize}
    \item \textbf{Physical Volume:} Because the singularity exists in equilibrium, the resulting core is a finite-volume entity. Information is not deleted but is encoded within the saturated geometry.
    \item \textbf{Quantum Consistency:} The maintenance of an equilibrium state ensures that the causal evolution of quantum states remains unitary, as the destructive nature of an infinite singularity is absent.
    \item \textbf{Structural Necessity:} The paradox is resolved by acknowledging that spacetime possesses a structural limit; once saturation is reached, equilibrium and information preservation become mechanical requirements.
\end{itemize}

\section{Observational Verification: S2 Orbit}
The equilibrium state of the core carries measurable consequences for stellar orbits near supermassive black holes. The SMP predicts a deviation in the Schwarzschild precession of the star S2, reflecting the localized equilibrium conditions near the horizon:
\begin{equation}
\Delta \phi_{SMP} = \Delta \phi_{GR} \left( \frac{8\pi}{8\pi + P_m} \right)
\end{equation}



\section{Conclusion}
The Stiffness Matching Principle allows us to redefine the singularity as a structural saturation limit that results in a state of physical equilibrium. By replacing mathematical infinities with this saturated equilibrium state, we resolve the Information Paradox and provide a stable, testable extension to General Relativity that preserves the fundamental laws of both gravity and quantum mechanics.

\begin{thebibliography}{9}
\bibitem{Skavberg2026} 
Skavberg, R. (2026). \textit{A Mechanical Route to G: Matching Dark-Matter Field Pressure to the Einstein-Hilbert Coupling}. Zenodo. DOI: 10.5281/zenodo.18393277.
\bibitem{BransDicke1961}
Brans, C., and Dicke, R. H. (1961). \textit{Mach's Principle and a Relativistic Theory of Gravitation}. Physical Review.
\bibitem{Hawking1976}
Hawking, S. W. (1976). \textit{Breakdown of predictability in gravitational collapse}. Physical Review D.
\bibitem{GRAVITY2020}
GRAVITY Collaboration. (2020). \textit{Detection of the Schwarzschild precession in the orbit of the star S2}. Astronomy \& Astrophysics.
\end{thebibliography}

\end{document}